\section{Relation to Contemporaneous Work}\label{app:comparison}

Qiushi Engine independently developed the structural proof presented here.
Wang's contemporaneous preprint establishes the same numerical lower
bound~\cite{wang2026twentyone}. We became aware of that work during release
preparation, after the research and original scientific report had been completed.

The comparison below concerns Wang's published binary proof and its
reproduction package, and the proof and research record documented here.
It distinguishes shared mathematics from differences in the finite
premises, deductions, and research contributions.

\paragraph{Shared mathematical foundation.}
Both proofs inherit certified restriction bounds and use occupation
capacities. Both obtain rank-one differences between high-rank first
factors and place them in a common affine row or column coset. Restriction
by an annihilator and quotienting the coefficient space describe the same
operation in dual conventions. D'Ambrosio provides an earlier precedent
for capacity arguments and tight restrictions~\cite{dambrosio2026}.

\begin{longtable}{@{}L{2.1cm}L{5.7cm}L{6.4cm}@{}}
\toprule
Aspect & Wang's binary proof & Qiushi Engine's structural proof\\
\midrule
\endfirsthead
\toprule Aspect & Wang's binary proof & Qiushi Engine's structural proof\\\midrule
\endhead
\bottomrule\endfoot
Finite premises
& Rank-one-span searches and backtracking strengthen a restriction table.
& A frozen earlier table supports eight additional quotient bounds,
certified by CNF/DRAT refutations and by closed Lean integer branch
certificates.\\
After coset reduction
& Classify 35 higher-rank configurations and test their rank-one
completions against the strengthened table.
& Combine affine-plane caps with the full-tensor rank sum to force
exactly the profile $(16,1,3)$ and equality at rank sum 27.\\
Decisive deduction
& Exhaustive capacity-constrained enumeration excludes every first-factor
profile; no second- or third-factor completion is needed.
& Equality gives $M_tP^{-1}M_s=\delta_{ts}M_t$. The tensor product formula
then couples all three factors and contradicts the required invertibility.\\
Verification
& Check restriction certificates and the complete profile enumeration.
& The original computational route checks the inherited certificate and
regenerated CNF/DRAT pairs. The Lean route checks integer branch
certificates and the structural deduction, with all dependencies
kernel-rechecked.\\
Further results
& Also establishes $R_{\mathbb F_3}(\langle2,3,3\rangle)=15$.
& The saturation lemma also constrains zero-excess rank-22 decompositions:
at most one first factor is invertible. With all 22 factors in the slot
nonzero, the profiles are $(17,5,0)$ and $(18,3,1)$; the zero-factor
alternative is retained in Section~\ref{sec:outlook}.\\
\end{longtable}

The two routes use the common coset geometry in different ways: Wang tests
first-factor capacities against a strengthened table, while the argument
here forces rank equality and derives cross-factor identities. The latter
also gives the zero-excess profile restrictions in Section~\ref{sec:outlook}.

\paragraph{Research development and open materials.}
Wang supplies a reproducible theorem package, including code and
certificates. Our report also documents how the structural proof
developed: rank-22 construction and deformation, quotient cores,
support-versus-completion tests, corrections, the return to the full
tensor, and recognition of saturation. Appendix~\ref{app:research-record}
integrates this development into an evidence-linked Meta-Trace, supported
by the thematic research archive. This is a contribution in the
documentation and study of autonomous mathematical research, alongside
the structural argument.

Readers interested in rank-one-span searches, symmetry-reduced profile
enumeration, or the ternary rectangular case will find complementary
methods in Wang's article~\cite{wang2026twentyone}. Read alongside the
saturation argument and Meta-Trace here, that work offers another
perspective on how finite-field constraints can lead to tensor-rank
lower bounds.
